\subsection{Claim-Level Evidence Governance}
\label{subsec:claim_governance}

The extraction model treated a study-level ``included'' decision and a
claim-level ``usable'' decision as different judgments. Inclusion established
that a study belonged to the primary technical corpus; it did not make every
statement or number in that study suitable for cross-study synthesis. Each
extracted claim was therefore retained with its study and report identifiers,
source locator, claim type, reported conditions, units, validation context,
and survey-use disposition.

Claims were separated into evidence statements, quantitative metric results,
and trade-off relationships. Quantitative use required an identifiable metric
definition and role, a preserved unit, an explicit or defensibly coded
measurement plane, an operating condition, and a source location. Optical
OSNR was not substituted for electrical SNR or ESNR without a source-provided
conversion model. Physical resolution, estimator accuracy or RMSE, CRB-type
bounds, and fiber spatial granularity were retained as distinct constructs.
Simulation, analytical, laboratory, prototype, and field evidence were also
kept as distinct validation contexts.

Claims were assigned one of four operational survey-use dispositions:
eligible for qualitative synthesis, eligible for quantitative synthesis,
context only, or quarantined conflict. Quantitative eligibility indicates that
a claim may be used numerically under its recorded conditions; it does not
mean that the claim is automatically comparable with every other quantitative
claim or eligible for statistical pooling. Context-only claims may inform
terminology, limitations, or interpretation but may not support primary
technical prevalence or performance conclusions. Quarantined claims are
excluded from quantitative synthesis until the conflict is resolved from an
authoritative source.

Source disagreements were preserved in a conflict register. Conflicting
values were not averaged, normalized by assumption, or replaced with a
preferred value solely because that value appeared more plausible. A study
could therefore remain survey-ready while one or more of its claims remained
restricted or quarantined.

\subsubsection{Current Phase-D Claim Dispositions}

The governed Phase-D evidence model contains 8,306 claims across 206 included
studies. Of these, 3,206 claims are eligible for qualitative synthesis, 4,997
are eligible for quantitative synthesis, 31 are restricted to contextual use,
and 72 are quarantined. These four dispositions reconcile exactly to the full
claim universe:

\begin{equation}
3{,}206 + 4{,}997 + 31 + 72 = 8{,}306.
\end{equation}

At study level, 175 studies are classified as survey-ready and 31 as
survey-ready with claim restrictions, accounting for all 206 included studies.
The conflict register contains 93 rows, including 2 rows classified as explicit
metric conflicts. Conflict-register rows and quarantined claims are different
counting units and must not be presented as interchangeable totals.

\begin{table}[htbp]
\caption{Phase-D claim-governance and study-use status.}
\label{tab:phase_d_claim_governance}
\centering
\begin{tabular}{lr}
\toprule
\textbf{Governance category} & \textbf{Count} \\
\midrule
All governed claims & 8,306 \\
Claims eligible for qualitative synthesis & 3,206 \\
Claims eligible for quantitative synthesis & 4,997 \\
Context-only claims & 31 \\
Quarantined claims & 72 \\
Survey-ready studies & 175 \\
Survey-ready studies with claim restrictions & 31 \\
Conflict-register rows & 93 \\
Explicit metric-conflict rows & 2 \\
\bottomrule
\end{tabular}
\end{table}

These classifications are survey-use controls rather than TQAF quality grades.
The workflow provenance documents investigator-supervised, AI-assisted, and
investigator-delegated adjudication, but independent full-corpus human verification
is not documented. The corpus must therefore not be described as independently
human-verified, double-reviewed, or independently adjudicated. Likewise, the
175/31 study statuses must not be interpreted as evidence that every extracted
claim has received independent human source verification.

% PHASE-E INTEGRATION MARKER:
% Phase-E TQAF results are supplied as a separate appraisal layer in
% 06_PHASE_E_TQAF_RESULTS_EN.tex.
% Do not relabel the Phase-D claim dispositions or 175/31 study-use statuses as
% quality scores.

% PHASE-F INTEGRATION MARKER:
% Phase-F synthesis denominators are derived from the eligible claim subset
% required by each analysis; see 07_PHASE_F_S1_S7_RESULTS_EN.tex. Do not use
% 8,306, 4,997, or 206 as a universal denominator for taxonomy, metric,
% trade-off, or application claims.
